\clearpage
\VerdierPhysicalPageAnchor{0236}%
\begingroup\small\setlength{\parskip}{0pt}\setlength{\emergencystretch}{1.5em}
\setlength{\abovedisplayskip}{4pt plus 2pt minus 2pt}%
\setlength{\belowdisplayskip}{4pt plus 2pt minus 2pt}%

\VerdierSemanticAnchor{ch4d-4-2-4}
\noindent\textbf{4.2.4.}\quad
$\operatorname{comp}(\mathcal A)$ 中 $\widetilde{\mathbb Z}$ 型滤过对象的态射，
以显然的方式诱导命题 (4.2.3) a) 所述的相应谱对象之间的态射。因此，命题
(4.2.3) 定义了从 $\operatorname{comp}(\mathcal A)$ 的滤过对象范畴到
取值于 $D(\mathcal A)$ 的谱对象范畴的函子。

\VerdierSemanticAnchor{ch4d-4-2-4-first-canonical}
设 $X$ 是 $\operatorname{comp}(\mathcal A)$ 的对象，并设
$X(q)\subset X$，$q\in\mathbb Z$，是它的第一典范滤过。
在这种情形下，命题 (4.2.3) 关联给滤过对象
$\bigl(X,X(q),q\in\mathbb Z\bigr)$ 的谱对象，正是与复形 $X$ 相联系的、
取值于 $D(\mathcal A)$ 的第一谱对象 (4.1.5)。

\VerdierSemanticAnchor{ch4d-4-2-4-nfold}
同样，设 $X$ 是 $\mathcal A$ 中的 $n$-重复形，$\int_n X=Y$ 是其关联单复形，
$m$ 是整数，$1\leq m\leq n$，并设
\[
0={}^{m}\!Y(-\infty)\subset\cdots\subset{}^{m}\!Y(q)
\subset\cdots\subset{}^{m}\!Y(+\infty)=Y
\]
是由 $X$ 的第 $m$ 个次数所定义的滤过 (4.1.6.9)。在这种情形下，命题
(4.2.3) 关联给滤过对象
$\bigl(Y,{}^{m}\!Y(q),q\in\mathbb Z\bigr)$ 的谱对象，正是由按第 $m$ 个次数
滤过的 $X$ 所定义的、取值于 $D(\mathcal A)$ 的谱对象 (4.1.8)。

\VerdierSemanticAnchor{ch4d-4-3}
\medskip\noindent\textbf{4.3. 第二典范谱对象。}

\VerdierSemanticAnchor{ch4d-4-3-1}
\noindent\textbf{4.3.1.}\quad
设 $X$ 是 $\operatorname{comp}(\mathcal A)$ 的对象。令：

\VerdierSemanticAnchor{ch4d-4-3-1-1}
\begin{equation}
\begin{aligned}
X(-\infty)&=0,\qquad X(+\infty)=X,\\
X(q)&=\cdots\longrightarrow X^{q-3}\longrightarrow X^{q-2}
\longrightarrow\operatorname{Ker}(d_X^{q-1})
\longrightarrow0\longrightarrow0\longrightarrow\cdots,
\qquad q\in\mathbb Z,
\end{aligned}
\tag{4.3.1.1}
\end{equation}
后一个复形的微分由 $X$ 的微分诱导。存在典范单态射：

\VerdierSemanticAnchor{ch4d-4-3-1-2}
\begin{equation}
0=X(-\infty)\subset\cdots\subset X(q)\subset X(q+1)
\subset\cdots\subset X(+\infty)=X
\tag{4.3.1.2}
\end{equation}
它们在 $X$ 上定义了一个 $\widetilde{\mathbb Z}$ 型滤过，称为 $X$ 的
\emph{第二典范滤过}。由命题 (4.2.3) 得到一个谱对象
$\bigl(X(p,q),\delta(p,q,r)\bigr)$，称为与复形 $X$ 相联系的、取值于
$D(\mathcal A)$ 的\emph{第二谱对象}。

\VerdierSemanticAnchor{ch4d-4-3-1-associated-functor}
$\operatorname{comp}(\mathcal A)$ 中的态射 $X\longrightarrow Y$ 保持第二典范滤过，
因而函子
\[
X\longmapsto\bigl(X(p,q),\delta(p,q,r)\bigr)
\]
是从 $\operatorname{comp}(\mathcal A)$ 到取值于 $D(\mathcal A)$ 的谱对象
范畴的函子。这个函子称为\emph{第二典范谱对象函子}。
\endgroup
